\documentclass[10pt]{article}

\usepackage[margin=0.75in]{geometry}
\usepackage{amsmath,amsthm,amssymb,bm}
\usepackage{xcolor}
\usepackage{cancel}
\usepackage{graphicx}
\usepackage{changepage}
\usepackage{circuitikz}
\usepackage{pgfplots}
\usepackage{physics}
\usepackage{hyperref}
\usepackage{siunitx}
\usepackage{minted}
\usepackage{fontspec}
\usepackage{relsize}
\usepackage{subfig}
\usepackage{mhchem}
\usepackage{todonotes}
\usepackage{biblatex}
\usepackage{multicol, multirow, booktabs}
\usepackage[breakable]{tcolorbox}
\usepackage[inline]{enumitem}
\patchcmd{\thebibliography}{\section*{\refname}}{}{}{}
\addbibresource{sources.bib}

\theoremstyle{definition}
\newtheorem{problem}{Problem}
\newtheorem{soln}{Solution}

\pgfplotsset{compat=newest}
\usetikzlibrary{lindenmayersystems}
\usetikzlibrary{arrows}
\usetikzlibrary{calc}
\usetikzlibrary{positioning, fit}
\usetikzlibrary{3d, perspective}

\definecolor{incolor}{HTML}{303F9F}
\definecolor{outcolor}{HTML}{D84315}
\definecolor{cellborder}{HTML}{CFCFCF}
\definecolor{cellbackground}{HTML}{F7F7F7}
\newcommand{\ui}{\hat{i}}
\newcommand{\uj}{\hat{j}}
\newcommand{\uk}{\hat{k}}
\newcommand{\ux}{\hat{x}}
\newcommand{\uy}{\hat{y}}
\newcommand{\uz}{\hat{z}}
\newcommand{\uv}[1]{\hat{\bm{{#1}}}}
\newcommand{\pr}[1]{{#1^\prime}}
\newcommand{\justif}[2]{&{#1}&\text{#2}}
\newcommand{\dul}[1]{\underline{\underline{#1}}}
\pgfdeclarelayer{background}  
\pgfsetlayers{background,main}
\AtBeginDocument{\RenewCommandCopy\qty\SI}

\makeatletter
\newcommand{\boxspacing}{\kern\kvtcb@left@rule\kern\kvtcb@boxsep}
\makeatother
\newcommand{\prompt}[4]{
    \ttfamily\llap{{\color{#2}[#3]:\hspace{3pt}#4}}\vspace{-\baselineskip}
}

\newcommand{\thevenin}[2]{
  \begin{center}
    \begin{circuitikz} \draw
      (0,0) -- (2,0) to[battery1, l_=$V_{Th}\eq#1$] (2,2) 
      to[resistor, l_=$R_{Th}\eq#2$] (0,2)
      ;
      \draw [o-] (-.07,2.079);
      \draw [o-] (-.07,0.079);
    \end{circuitikz}
  \end{center}
}

\newcommand{\norton}[2]{
  \begin{center}
    \begin{circuitikz} \draw
      (0,0) -- (3,0) to[american current source, l_=$I_{N}\eq#1$] (3,2) -- (0,2) (2,0)
      to[resistor, l=$R_{N}\eq#2$] (2,2)
      ;
      \draw [o-] (-.07,2.079);
      \draw [o-] (-.07,0.079);
    \end{circuitikz}
  \end{center}
}

\newcommand{\highlight}[1]{\colorbox{yellow}{$\displaystyle #1$}}

\newcommand{\ti}[1]{\widetilde{#1}}

\newfontface{\Kaufmann}{Kaufmann}
\DeclareTextFontCommand{\kf}{\Kaufmann}
\newcommand{\scriptr}{\fontsize{12pt}{12pt}\kf{r}}

\newfontface{\KaufmannB}{Kaufmann Bold}
\DeclareTextFontCommand{\kfb}{\KaufmannB}
\newcommand{\bscriptr}{\fontsize{12pt}{12pt}\kfb{r}}

\newcommand{\bv}[1]{\vec{#1}}

\title{Physics 4605H: Assignment IV}
\author{Jeremy Favro (0805980) \\ Trent University, Peterborough, ON, Canada}
\date{\today}

\begin{document}
\maketitle

% PROBLEM 1

\begin{problem}
Find the eigenvalues and eigenstates of $\hat{\sigma}_y$.
\end{problem}
\begin{soln}
  We solve the equation $\hat{\sigma}_y\bv{v}=\lambda\bv{v}$ (the eigenvalue equation) by taking
  $$\abs{\hat{\sigma}_y-\lambda I}=0\implies \begin{vmatrix}
      -\lambda & -i       \\
      i        & -\lambda
    \end{vmatrix}=0\implies \lambda^2-1=0\implies \lambda=\pm 1$$. We then evaluate the eigenvalue equation for each $\lambda$,
  $$\begin{pmatrix}
      0 & -i \\
      i & 0
    \end{pmatrix}\begin{pmatrix}
      a \\b
    \end{pmatrix}=\begin{pmatrix}
      a \\
      b
    \end{pmatrix}\implies -ib=a,\,ia=b.$$
  This doesn't tell us much on its own, however we also have that (or rather, we want that) the eigenvectors should be normalized, hence we choose
  $$a=\frac{1}{\sqrt{2}}\implies \bv{v}_1=\frac{1}{\sqrt{2}}\begin{pmatrix}
      1 \\i
    \end{pmatrix}.$$
  We follow the same process for $\lambda=-1$,
  $$\begin{pmatrix}
      0 & -i \\
      i & 0
    \end{pmatrix}\begin{pmatrix}
      -a \\-b
    \end{pmatrix}=\begin{pmatrix}
      a \\
      b
    \end{pmatrix}\implies ib=a,\,-ia=b$$
  and again enforce normalization to obtain that, with $a=1/\sqrt{2}$,
  $$\bv{v}_{-1}=\begin{pmatrix}
      1 \\
      -i
    \end{pmatrix}$$
\end{soln}

% PROBLEM 2
\begin{problem}
For each of the scenarios below, find the corresponding density operator, both (i) in
Dirac notation and (ii) in matrix notation. Also, in each case (iii) evaluate the trace.
\begin{enumerate}[label=(\alph*)]
  \item A single-qubit system which is in the pure quantum state $\ket{0}$.
  \item A single-qubit system which is in the pure quantum state $\ket{+}\equiv \left(\ket{0}+\ket{1}\right)/\sqrt{2}$.
  \item A single-qubit system is in a mixed state in which there is a probability $1/4$ of the state
        $\ket{0}$ and probability $3/4$ of the state $\ket{+}$.
\end{enumerate}
\end{problem}
\begin{soln}~
  \begin{enumerate}[label=(\alph*)]
    \item We defined the density operator to be $\hat{\rho}=\ket{0}\bra{0}$. In matrix notation this is given as
          $$\begin{pmatrix}
              1 \\0
            \end{pmatrix}\begin{pmatrix}
              1 & 0
            \end{pmatrix}=\begin{pmatrix}
              1 & 0 \\
              0 & 0
            \end{pmatrix}.$$
          Here $\tr(\rho)=1+0=1$.
    \item Again,
          $$\hat{\rho}=\frac{1}{2}\left(\ket{0}+\ket{1}\right)\left(\bra{0}+\bra{1}\right)
            =\frac{1}{2}\left[\ket{0}\bra{0}+\ket{0}\bra{1}+\ket{1}\bra{0}+\ket{1}\bra{1}\right]
          $$ and
          $$
            \dul{\rho}=\frac{1}{2}\begin{pmatrix}
              1 \\1
            \end{pmatrix}\begin{pmatrix}
              1 & 1
            \end{pmatrix}
            =\frac{1}{2}\begin{pmatrix}
              1 & 1 \\
              1 & 1
            \end{pmatrix}.
          $$
          Here $\tr(\rho)=\frac{1}{2}\left(1+1\right)=1$.
    \item The density matrix here will be a weighted average of the pure-state density operators for each state, which we found in the previous question, and with
          weights given by the probabilities we are given here. Hence,
          $$\hat{\rho}=\frac{1}{4}\hat{\rho_{\ket{0}}}+\frac{3}{4}\hat{\rho}_{\ket{+}}
            =\frac{1}{4}\ket{0}\bra{0}+\frac{3}{8}\left[\ket{0}\bra{0}+\ket{0}\bra{1}+\ket{1}\bra{0}+\ket{1}\bra{1}\right]
          $$
          which is, in matrix form,
          $$\frac{1}{4}\begin{pmatrix}
              1 & 0 \\
              0 & 0
            \end{pmatrix}+\frac{3}{8}\begin{pmatrix}
              1 & 1 \\
              1 & 1
            \end{pmatrix}=\begin{pmatrix}
              5/8 & 3/8 \\
              3/8 & 3/8
            \end{pmatrix}.$$
          Here $\tr(\rho)=\left(5/8+3/8\right)=1$.
  \end{enumerate}
\end{soln}

% PROBLEM 3
\begin{problem}
Consider a system consisting of two qubits, $A$ and $B$. Qubit $A$ is in the state described
in (2a). Qubit $B$ is in the state described by (2b).
\begin{enumerate}[label=(\alph*)]
  \item What is the density matrix $\rho_{AB}$ of the combined system? Give the result both in Dirac
        notation and in matrix notation (using the basis $\ket{00},\ket{01},\ket{10},\ket{11}$).
  \item Calculate the partial trace $\rho_A=\tr_B(\rho_{AB})$
\end{enumerate}
\end{problem}
\begin{soln}~
  \begin{enumerate}[label=(\alph*)]
    \item In Dirac notation
          \begin{align*}
            \rho_{AB} & =\rho_A\otimes \rho_B                                                                                                                                                    \\
                      & =(\ket{0}\bra{0})\otimes\frac{1}{2}\left[\ket{0}\bra{0}+\ket{0}\bra{1}+\ket{1}\bra{0}+\ket{1}\bra{1}\right]                                                              \\
                      & =\frac{1}{2}\left[\ket{0}\bra{0}\otimes\ket{0}\bra{0}+\ket{0}\bra{0}\otimes\ket{0}\bra{1}+\ket{0}\bra{0}\otimes\ket{1}\bra{0}+\ket{0}\bra{0}\otimes\ket{1}\bra{1}\right] \\
                      & =\frac{1}{2}\left[\ket{00}\bra{00}+\ket{00}\bra{01}+\ket{01}\bra{00}+\ket{01}\bra{01}\right]                                                                             \\
          \end{align*}
          and in matrix notation,
          $$\rho_{AB}=\rho_A\otimes \rho_B
            =\frac{1}{2}\begin{pmatrix}
              1 & 0 \\
              0 & 0
            \end{pmatrix}\otimes\begin{pmatrix}
              1 & 1 \\
              1 & 1
            \end{pmatrix}=\frac{1}{2}\begin{pmatrix}
              1 & 1 & 0 & 0 \\
              1 & 1 & 0 & 0 \\
              0 & 0 & 0 & 0 \\
              0 & 0 & 0 & 0
            \end{pmatrix}
          $$
    \item This I'll do in Dirac notation because I can't make sense of my notes (or the textbook) on this for how it would work in matrix notation. I think I'll remember to ask about it at some point.
          The general expression for \emph{tracing out} $B$ to obtain $\rho_A$ is
          $$\rho_A=\tr_B(\rho_{AB})=\sum_i\bra{B_i}\rho_{AB}\ket{B_i}$$
          which I think here is
          \begin{align*}
            \rho_A & =\sum_i\bra{B_i}\rho_{AB}\ket{B_i}                                                                                                                                                                  \\
                   & =\frac{1}{2}\left(\bra{0}\left[\ket{0}\bra{0}\otimes\ket{0}\bra{0}+\ket{0}\bra{0}\otimes\ket{0}\bra{1}+\ket{0}\bra{0}\otimes\ket{1}\bra{0}+\ket{0}\bra{0}\otimes\ket{1}\bra{1}\right]\ket{0}\right. \\
                   & \quad+\left.\bra{1}\left[\ket{0}\bra{0}\otimes\ket{0}\bra{0}+\ket{0}\bra{0}\otimes\ket{0}\bra{1}+\ket{0}\bra{0}\otimes\ket{1}\bra{0}+\ket{0}\bra{0}\otimes\ket{1}\bra{1}\right]\ket{1}\right)
            % &=\frac{1}{2}\left(\bra{00}\left[\ket{00}\bra{00}+\ket{00}\bra{01}+\ket{01}\bra{00}+\ket{01}\bra{01}\right]\ket{00}\right.\\
            % &\quad+\bra{01}\left[\ket{00}\bra{00}+\ket{00}\bra{01}+\ket{01}\bra{00}+\ket{01}\bra{01}\right]\ket{01}\\
            % &\quad+\bra{10}\left[\ket{00}\bra{00}+\ket{00}\bra{01}+\ket{01}\bra{00}+\ket{01}\bra{01}\right]\ket{10}\\
            % &\quad+\left.\bra{11}\left[\ket{00}\bra{00}+\ket{00}\bra{01}+\ket{01}\bra{00}+\ket{01}\bra{01}\right]\ket{11}\right)
          \end{align*}
          This is where I got tripped up. After doing some reading I've realised that when I act on the states in the square brackets I can only act on the states in $B$ as the states
          I'm acting with are in $B$. I'll dissassemble the tensor product to make this more clear,
          \begin{align*}
            &=\frac{1}{2}(\ket{0}\bra{0})\otimes
            \left(\bra{0}\left[\ket{0}\bra{0}+\ket{0}\bra{1}+\ket{1}\bra{0}+\ket{1}\bra{1}\right]\ket{0}+
             \bra{1}\left[\ket{0}\bra{0}+\ket{0}\bra{1}+\ket{1}\bra{0}+\ket{1}\bra{1}\right]\ket{1}\right)\\
             &=\frac{1}{2}(\ket{0}\bra{0})\otimes
            \left(\left[\bra{0}\ket{0}\bra{0}\ket{0}+\bra{0}\ket{0}\bra{1}\ket{0}+\bra{0}\ket{1}\bra{0}\ket{0}+\bra{0}\ket{1}\bra{1}\ket{0}\right]+
             \left[\bra{1}\ket{0}\bra{0}\ket{1}+\bra{1}\ket{0}\bra{1}\ket{1}+\bra{1}\ket{1}\bra{0}\ket{1}+\bra{1}\ket{1}\bra{1}\ket{1}\right]\right)\\
             &=\frac{1}{2}(\ket{0}\bra{0})\otimes
            \left(\left[1+0+0+0\right]+
             \left[0+0+0+1\right]\right)=\frac{2}{2}\begin{pmatrix}
              1&0\\
              0&0
             \end{pmatrix}
          \end{align*}
          which successfully recovers our original matrix.

          % which obviously isn't the same as the $\rho_A$ we had previously. I have no idea where I'm going wrong here as combined $\rho_{AB}$ gives the correct(? I think)
          % matrix. Maybe I am misunderstanding what $\left\{\ket{B_i}\right\}$ is. I am under the impression that it is the set of basis vectors for the $B$ state,
          % which is why I can't make sense of how it would work in matrix notation as that would be a 4x4 operating on a 2x1. I think in Dirac that the basis would be
          % $\left\{\ket{0},\ket{1}\right\}$ as I tried above, but maybe it's $\left\{\ket{+},\ket{-}\right\}$? I'll try that,
          % \begin{align*}
          %   \rho_A & =\sum_i\bra{B_i}\rho_{AB}\ket{B_i}                                                                                                                                                                                      \\
          %          & =\frac{1}{4}\left((\bra{0}+\bra{1})\left[\ket{0}\bra{0}\otimes\ket{0}\bra{0}+\ket{0}\bra{0}\otimes\ket{0}\bra{1}+\ket{0}\bra{0}\otimes\ket{1}\bra{0}+\ket{0}\bra{0}\otimes\ket{1}\bra{1}\right](\ket{0}+\ket{1})\right. \\
          %          & \quad+\left.(\bra{0}-\bra{1})\left[\ket{0}\bra{0}\otimes\ket{0}\bra{0}+\ket{0}\bra{0}\otimes\ket{0}\bra{1}+\ket{0}\bra{0}\otimes\ket{1}\bra{0}+\ket{0}\bra{0}\otimes\ket{1}\bra{1}\right](\ket{0}-\ket{1})\right)
          % \end{align*}
          % I've just realized this won't work as I'll jsut get twice the original matrix, times $1/4$ which will give me back the original because the $\ket{1}$ cases are all zero again and so I just have two of the $\ket{0}$ ones.
          % So no luck.
  \end{enumerate}
\end{soln}

% PROBLEM 4
\begin{problem}
Determine the Bloch sphere angles $\theta$ and $\phi$ for the eigenstates of each Pauli matrix.
\end{problem}
\begin{soln}
  The eigenvectors of the Pauli matrices are
  \begin{align*}
    \psi_{x+} & = \frac{1}{\sqrt{2}} \begin{pmatrix} 1 \\  1 \end{pmatrix}, &
    \psi_{x-} & = \frac{1}{\sqrt{2}} \begin{pmatrix} 1 \\ -1 \end{pmatrix},   \\
    \psi_{y+} & = \frac{1}{\sqrt{2}} \begin{pmatrix} 1 \\  i \end{pmatrix}, &
    \psi_{y-} & = \frac{1}{\sqrt{2}} \begin{pmatrix} 1 \\ -i \end{pmatrix},   \\
    \psi_{z+} & =                  \begin{pmatrix} 1 \\  0 \end{pmatrix},   &
    \psi_{z-} & =                  \begin{pmatrix} 0 \\  1 \end{pmatrix}.
  \end{align*}
  We express state vectors on the Bloch sphere as
  $$\ket{\psi}=\cos(\theta/2)\ket{0}+e^{i\phi}\sin(\theta/2)\ket{1}.$$
  We can then consider which values of $\theta$ and $\phi$ will give us the desired states from above. This is not hugely complex so I will just tabulate my results:

  \begin{table}[ht!]
    \centering%
    \begin{tabular}{lll}
      \toprule
      State       & $\theta$ & $\phi$   \\
      \midrule
      $\psi_{x+}$ & $\pi/2$  & $0$      \\
      $\psi_{x-}$ & $\pi/2$  & $\pi$    \\
      $\psi_{y+}$ & $\pi/2$  & $\pi/2$  \\
      $\psi_{y-}$ & $\pi/2$  & $3\pi/2$ \\
      $\psi_{z+}$ & $0$      & Free     \\
      $\psi_{z-}$ & $\pi$    & $0$      \\
      \bottomrule
    \end{tabular}
  \end{table}
\end{soln}
\end{document}